%% ============================================================
%% 参考文献（自动另起一页；“前引后列”：每条都须在正文中被 \cite 引用，
%% 并按引用先后顺序排列；著录格式按 GB/T 7714-2015）
%% ============================================================
\begin{thebibliography}{99}

%% 期刊文献 [J]：主要责任者. 题名[J]. 期刊名，年，卷(期):页码.
\bibitem{ref1} 张三，李四，王五，等. 凸轮分度机构动力学建模与仿真研究[J]. 机械工程学报，2024，60(3):45-52.

%% 学位论文 [D]：主要责任者. 题名[D]. 大学所在城市:大学名称，出版年.
\bibitem{ref2} 赵六. 高速凸轮机构动态特性分析[D]. 天津:天津理工大学，2023.

%% 普通图书 [M]：主要责任者. 书名[M]. 版本项. 出版地:出版者，出版年:引文页码.
\bibitem{ref3} 孙七. 机械原理[M]. 3版. 北京:机械工业出版社，2021:120-135.

%% 电子资源 [EB/OL]：必须列出资料所在的具体页面、引用日期与访问路径
\bibitem{ref4} 周八. 工业凸轮机构设计手册[EB/OL]. (2023-05-12)[2026-03-01]. https://www.example.com/cam-handbook.

%% 外文期刊（作者超过三位列前三位，后加 et al.）
\bibitem{ref5} Smith J, Brown K, Davis M, et al. Dynamic analysis of indexing cam mechanisms[J]. Mechanism and Machine Theory, 2023, 180: 105-118.

\end{thebibliography}
